\documentclass[11pt, a4paper]{article}

\usepackage[T1]{fontenc}
\usepackage[utf8]{inputenc}
\usepackage{tgpagella}
\usepackage{tgcursor}
\usepackage{microtype}
\usepackage[spanish, es-noshorthands]{babel}
\usepackage[table, xcdraw]{xcolor}
\usepackage{booktabs}
\usepackage{tabularx}
\usepackage{array}
\usepackage{amsmath, amssymb}
\usepackage[most]{tcolorbox}
\usepackage[margin=2cm]{geometry}
\usepackage{fancyhdr}

\pagestyle{fancy}
\fancyhf{}
\setlength{\headheight}{16pt}
\lhead{\textbf{IMT-342}}
\chead{Checklist Operativo de Seguridad y Autorización (Gate 3)}
\rhead{Semana 12 - Sesión 1}
\lfoot{Protocolo GO / NO-GO}
\rfoot{Página \thepage}
\renewcommand{\headrulewidth}{0.4pt}
\renewcommand{\footrulewidth}{0.4pt}

\definecolor{ucbblue}{RGB}{0, 51, 102}

\begin{document}

\begin{center}
    \Large\textbf{\textcolor{red!80!black}{Protocolo y Checklist de Seguridad para Pruebas en Hardware}}[cite: 15]
\end{center}

\begin{tcolorbox}[colback=red!5, colframe=red!80!black, title=\textbf{Estado Inicial de la Sesión: NO-GO para Movimiento Físico} {[cite: 15]}, left=2mm, right=2mm, top=1mm, bottom=1mm]
    \small
    Ninguna prueba con movimiento físico procede hasta que el operador autorizado y el responsable del laboratorio completen la verificación externa y declaren expresamente \texttt{GO}[cite: 15]. Ante cualquier campo crítico pendiente, duda sobre marcos/objetivos o falta de supervisión formal, la decisión obligatoria es \textbf{\texttt{NO-GO}}[cite: 15].
\end{tcolorbox}

\section*{1. Distinción de Niveles de Prueba y Evidencia[cite: 15]}
\begin{center}
\footnotesize
\renewcommand{\arraystretch}{1.35}
\begin{tabularx}{\textwidth}{|>{\hsize=0.8\hsize\raggedright\arraybackslash}X|>{\hsize=0.9\hsize\raggedright\arraybackslash}X|>{\hsize=1.15\hsize\raggedright\arraybackslash}X|>{\hsize=1.15\hsize\raggedright\arraybackslash}X|}
\hline
\rowcolor{gray!20}
\textbf{Nivel de Prueba}[cite: 15] & \textbf{Ejemplo}[cite: 15] & \textbf{Lo que Demuestra}[cite: 15] & \textbf{Lo que NO Demuestra}[cite: 15] \\ \hline
Simulación offline[cite: 15] & FK/MoveIt/RViz sin controlador[cite: 15] & Coherencia de software/modelo bajo supuestos[cite: 15] & Estado real del robot ni seguridad física[cite: 15] \\ \hline
Conectividad sin movimiento[cite: 15] & Respuesta del controlador bajo interfaz autorizada[cite: 15] & Presencia de comunicación determinada[cite: 15] & Control de trayectoria ni autorización de movimiento[cite: 15] \\ \hline
Estado observado[cite: 15] & \texttt{joint\_states} verificablemente originados en ABB[cite: 15] & Adquisición de datos y trazabilidad de lectura[cite: 15] & Aceptación de comando de movimiento[cite: 15] \\ \hline
Movimiento supervisado[cite: 15] & Objetivo ejecutado y medido en zona protegida[cite: 15] & Evidencia física acotada del Gate 3[cite: 15] & Seguridad certificada o validez en cualquier condición[cite: 15] \\ \hline
\end{tabularx}
\end{center}

\section*{2. Checklist de Puertas de Seguridad Previas a Gate 3[cite: 15]}
\begin{center}
\footnotesize
\renewcommand{\arraystretch}{1.45}
\begin{tabularx}{\textwidth}{|>{\hsize=1.25\hsize\raggedright\arraybackslash}X|>{\hsize=1.05\hsize\raggedright\arraybackslash}X|>{\hsize=0.85\hsize\raggedright\arraybackslash}X|c|>{\hsize=0.85\hsize\raggedright\arraybackslash}X|}
\hline
\rowcolor{ucbblue!15}
\textbf{Condición Crítica}[cite: 15] & \textbf{Cómo se Verifica}[cite: 15] & \textbf{Autoridad que Confirma}[cite: 15] & \textbf{Estado}[cite: 15] & \textbf{Evidencia / Fecha}[cite: 15] \\ \hline
1. Identidad / controlador / manual correcto[cite: 15] & Inspección e inventario[cite: 15] & Responsable técnico[cite: 15] & Pendiente & \\ \hline
2. Autorización de acceso al robot[cite: 15] & Registro institucional[cite: 15] & Laboratorio[cite: 15] & Pendiente & \\ \hline
3. Operador habilitado y supervisor presente[cite: 15] & Asignación formal[cite: 15] & Laboratorio[cite: 15] & Pendiente & \\ \hline
4. Modo operativo permitido[cite: 15] & Protocolo / manual concreto[cite: 15] & Operador autorizado[cite: 15] & Pendiente & \\ \hline
5. Zona de exclusión y despeje físico[cite: 15] & Inspección situacional[cite: 15] & Operador autorizado[cite: 15] & Pendiente & \\ \hline
6. Parada de emergencia (E-stop) y parada normal verificadas[cite: 15] & Procedimiento institucional[cite: 15] & Operador autorizado[cite: 15] & Pendiente & \\ \hline
7. Trayectoria y objetivos inspeccionados[cite: 15] & Validación offline y revisión[cite: 15] & Operador / Docente[cite: 15] & Pendiente & \\ \hline
8. Límites, velocidad reducida, tool/TCP y base[cite: 15] & Modelo / doc. / revisión[cite: 15] & Operador autorizado[cite: 15] & Pendiente & \\ \hline
9. Riesgo de objeto/pieza/gripper minimizado (sin pieza peligrosa en 1er ensayo)[cite: 15] & Configuración y evaluación local[cite: 15] & Operador autorizado[cite: 15] & Pendiente & \\ \hline
10. Registro de estado y mecanismos de cancelación/aborto[cite: 15] & Protocolo aprobado[cite: 15] & Operador autorizado[cite: 15] & Pendiente & \\ \hline
11. Condiciones para abortar y recuperar[cite: 15] & Procedimiento aprobado[cite: 15] & Laboratorio[cite: 15] & Pendiente & \\ \hline
\end{tabularx}
\end{center}

\vspace{0.3cm}
\begin{tcolorbox}[colback=white, colframe=black, title=\textbf{Registro de Decisión para Prueba Específica} {[cite: 15]}, left=2mm, right=2mm, top=2mm, bottom=2mm]
    \small
    \textbf{Decisión Operativa[cite: 15]:} \qquad $\square$ \textbf{\texttt{GO}} \qquad\qquad $\square$ \textbf{\texttt{NO-GO}} \\[0.3cm]
    \textbf{Tipo de prueba autorizada (alcance exacto)[cite: 15]:} \rule{9.5cm}{0.15mm} \\[0.3cm]
    \textbf{Firma e identidad del operador autorizado[cite: 15]:} \rule{9.2cm}{0.15mm} \\[0.3cm]
    \textbf{Fecha, hora y referencia al protocolo local[cite: 15]:} \rule{9.2cm}{0.15mm}
    \vspace{0.15cm}
    \par\scriptsize\textit{Nota importante: Esta plantilla no certifica ni sustituye las normas oficiales de ABB ni los reglamentos locales de seguridad del laboratorio. Prohibido puentear resguardos o improvisar secuencias de arranque[cite: 15].}
\end{tcolorbox}

\end{document}
